\documentclass[11pt,a4paper]{report}

\usepackage[margin=2.3cm]{geometry}
\usepackage[english]{babel}
\usepackage[T1]{fontenc}
\usepackage[latin1]{inputenc}
\usepackage{graphicx}
\usepackage{titlesec, blindtext, color}
\usepackage{fancyhdr}
\usepackage{hyperref}
\usepackage{float}
\usepackage{lastpage}


%----------------------------------------------------------------------------------------
%   Title Style Chapter
%----------------------------------------------------------------------------------------
\definecolor{gray75}{gray}{0.65}
\newcommand{\hsp}{\hspace{15pt}}
\titleformat{\chapter}[hang]{\Huge\bfseries}{\thechapter\hsp\textcolor{gray75}{|}\hsp}{0pt}{\Huge\bfseries}


%----------------------------------------------------------------------------------------
%   Link Style in ToC
%----------------------------------------------------------------------------------------
\hypersetup{
    colorlinks,
    citecolor=black,
    filecolor=black,
    linkcolor=black,
    urlcolor=black
}


%----------------------------------------------------------------------------------------
%   Header/Footer
%----------------------------------------------------------------------------------------
\pagestyle{fancy} 
%\lhead{\sectionmark}
%\chead{}
%\rhead{}
%\lfoot{}
%%\cfoot{\thepage/\pageref{LastPage}}
%\cfoot{\thepage}
%\rfoot{}
%\renewcommand{\headrulewidth}{0.3pt}
%%\renewcommand{\footrulewidth}{0.3pt}

\renewcommand{\chaptermark}[1]{\markboth{#1}{}}
\fancyhf{}
\fancyhead[C]{\leftmark}
\fancyfoot[C]{\thepage}

% New Page for each section
%\newcommand{\sectionbreak}{\clearpage}

% Renaming of ToC
\addto\captionsenglish{% Replace "english" with the language you use
  \renewcommand{\contentsname}%
    {Table of Contents}%
}

%----------------------------------------------------------------------------------------
%   START DOCUMENT 
%----------------------------------------------------------------------------------------
\begin{document}

%----------------------------------------------------------------------------------------
%	TITLEPAGE
%----------------------------------------------------------------------------------------
\begin{titlepage}

\newcommand{\HRule}{\rule{\linewidth}{0.3mm}} % Defines a new command for the horizontal lines, change thickness here

 \vspace*{0.3cm}
 
\center % Center everything on the page

%----------------------------------------------------------------------------------------
%	LOGO SECTION
%----------------------------------------------------------------------------------------
\includegraphics[scale=0.33]{../logo.png}\\[0.4cm] 

%----------------------------------------------------------------------------------------
%	HEADING SECTION
%----------------------------------------------------------------------------------------
\textsc{\Large \textbf{Politecnico di Milano}\\MSc Computer Science and Engineering}\\[1.7cm] \textsc{\Large \textbf{Software Engineering 2}\\[0.1cm]Academic Year 2016-2017}\\[1cm] 
%\textsc{\LARGE Minor Heading}\\[0.5cm] 

%----------------------------------------------------------------------------------------
%	TITLE SECTION
%----------------------------------------------------------------------------------------

\HRule \\[0.4cm]
{ \huge \bfseries Requirements Analysis and\\Specification Document\\[0.5cm]\textit{PowerEnJoy}}\\[0.4cm] % Title of document
\HRule \\[1.5cm]
 
%----------------------------------------------------------------------------------------
%	AUTHOR SECTION
%----------------------------------------------------------------------------------------
\Large \emph{Author:}\\
\LARGE{
\begin{center}
\begin{tabular}{ r  l }
	Melloni Giulio & 876279\\
	Renzi Marco & 878269\\
	Testa Filippo & 875456 \\
\end{tabular}
\end{center}
}
\vspace*{0.8cm}
\Large{\emph{Reference Professor:} \\
{\textsc{Mottola} Luca} 
}

%----------------------------------------------------------------------------------------
%	DATE SECTION
%----------------------------------------------------------------------------------------
 \vspace*{1cm}
{\LARGE {\it Release Date: October 25$^{th}$, 2016\\Version 0.1}}
%vfill % Fill the rest of the page with whitespace

\end{titlepage}


%----------------------------------------------------------------------------------------
%   TOC
%----------------------------------------------------------------------------------------
\tableofcontents

%----------------------------------------------------------------------------------------
%  BEGIN TEXT
%----------------------------------------------------------------------------------------
\chapter{Introduction}
\section{Description of the problem}
%\section{Purpose}
% is the same of section just above??
\section{Scope}
%identifies system-to-be and in its application domain
% GOAL HERE?
\section{Definitions, acronyms and abbreviations}
%glossary and similar...
\section{Reference Documents}
% ???
\section{Overview}
% summary description of the section of the document
\chapter{Overall Description}
\section{Product perspective}
% external interfaces (user, hw+sw), operations and site adaptions...
\section{Product functions}
% major functions  (GOAL HERE?)
\section{User characteristics}
\section{Constraints}
% any kind of limitations in the development
\section{Assumptions}
\subsection{Text assumptions}
% TODO
\subsection{Domain assumptions}
% TODO
\chapter{Specific Requirements}
\section{External Interface Requirements}
\subsection{User Interfaces}
\subsection{Hardware Interfaces}
\subsection{Software Interfaces}
\subsection{Communication Interfaces}
\section{Functional Requirements}
% deriving requirements from goals'
\section{Non Functional Requirements}
\chapter{Scenarios}
\chapter{Modeling the System}
\section{UML Models}
\subsection{Use Cases}
\subsection{Class Diagram}
\subsection{Sequence Diagram}
\subsection{Statechart Diagram}
\subsection{Activity Diagram}
\section{Alloy}
\subsection{Source Code}
%\subsection{}
\chapter{RASD Preparation}
\section{Tools}
\section{Timing}


%\begin{figure}[H]
%\centering
%\fbox{\includegraphics[width=\textwidth,keepaspectratio]{../file.jpg}}
%\caption{Caption}
%\end{figure}

%----------------------------------------------------------------------------------------
%   END DOCUMENT
%----------------------------------------------------------------------------------------
\end{document}